\documentclass[10pt]{article}

% Specify the margins and text size.
\setlength{\textwidth}{6.4in}
\setlength{\textheight}{9.8in}
\setlength{\oddsidemargin}{0pt}
\setlength{\evensidemargin}{0pt}
\setlength{\topmargin}{0pt}
\setlength{\hoffset}{.05in}
\setlength{\voffset}{-1in}

\setlength{\parskip}{5pt}
\setlength{\parindent}{0pt}

% Load some fonts and symbol packages
\usepackage{latexsym}
\usepackage{pifont}       % contains 'star' symbol for counterinsurgency handbook title
%\usepackage{yfonts} 
\usepackage{amsmath}
\usepackage{amsfonts}

\usepackage{graphicx}     % actually, this is loaded by pstricks
\usepackage[T1]{fontenc}
\usepackage{ifthen}
\usepackage{pstricks,pst-grad,pst-text,pst-node,multido,pst-plot,calc,pst-3dplot}
%\usepackage[all]{xy}
%\usepackage{animate}

% The hyperref package inserts links into the pdf.
\definecolor{MyLinkColor}{rgb}{.1,.2,1}
\definecolor{MyCiteColor}{rgb}{.1,1,.2}
\definecolor{MyURLColor}{rgb}{.4,.4,.4}
\usepackage[backref=true,pagebackref=false,hyperindex,colorlinks=true,
  linkcolor=MyLinkColor,urlcolor=MyURLColor]{hyperref}


% The tweaklist package is something I found on the web.  It provides a simple interface
% for making changes to spacing used in the itemize and enumerate environments.  Comment
% this out if you don't care to use tweaklist.
\usepackage{tweaklist}
\renewcommand{\itemhook}{\setlength{\parskip}{2pt}\setlength{\parsep}%
{1pt}\setlength{\topsep}{0pt}\setlength{\itemsep}{0pt}}

\newcommand{\U}{\underline{\hspace{5pt}}}

\usepackage{listings}
\newcommand{\Z}{\hphantom{0}}

\begin{document}
\pagestyle{empty}
\lstset{language=R, showspaces=false, showstringspaces=false}
\href{http://www.su.edu}{\includegraphics[height=1.75cm]{sulogo.eps}}
\vspace{-1.79cm}

{{\ }\hfill\small
\begin{tabular}{cl}
& Math 207\\
& Introduction to Statistics\\
\end{tabular}
}

\setlength{\baselineskip}{1.05\baselineskip}
\medskip

\begin{center}
\textbf{\large  Expected Values and Standard Errors for Sums}
\end{center}

\newcommand{\HH}{\hspace{20pt}\hphantom{a) } }
%
\newcommand{\twobox}[2]{\begin{pspicture}(-0.3,0.2)(2,1.3)
\psframe(0,0)(0.7,0.7)\rput(0.35,0.35){#1}
\psframe(1,0)(1.7,0.7)\rput(1.35,0.35){#2}
\psline(-0.3,0.9)(-0.3,-0.2)(2,-0.2)(2,0.9)(-0.3,0.9)
\end{pspicture}}
%
\newcommand{\threebox}[3]{\begin{pspicture}(0.2,0.2)(2.5,1.3)
\psframe(0,0)(0.7,0.7)\rput(0.35,0.35){#1}
\psframe(1,0)(1.7,0.7)\rput(1.35,0.35){#2}
\psframe(2,0)(2.7,0.7)\rput(2.35,0.35){#3}
\psline(-0.3,0.9)(-0.3,-0.2)(3,-0.2)(3,0.9)(-0.3,0.9)
\end{pspicture}}


1. Fill in the blanks. Compute the $z$-scores for the range then use
the normal curve to compute the probability of getting a sample sum
in that range.
\vspace{-5pt}

\begin{center}
{\setlength{\tabcolsep}{8pt}
\begin{tabular}{cccccccc}
\textbf{Box} & $\mbox{mean}_{\mbox{\footnotesize box}}$ & 
   $\mbox{SD}_{\mbox{\footnotesize box}}$ & $n$ & $\mbox{EV}_{\mbox{\footnotesize sum}}$ &
   $\mbox{SE}_{\mbox{\footnotesize sum}}$ & \textbf{range} & \textbf{probability}\\[3pt]\hline
\twobox{4}{6}      & 5 & 1 & 16     & 80 & 4 & 76 to 84 & 68\%\\[8pt]
\twobox{4}{6}      & & & 16     & & & 80 to 88 &\\[8pt]
\twobox{4}{6}      & & & 16     & & & over 84 &\\[8pt]
\twobox{4}{6}      & & & 100    & & & 490 to 510 &\\[8pt]
\twobox{4}{6}      & & & 100    & & & 485 to 515 &\\[8pt]
\twobox{4}{6}      & & & 100    & & & over 520 &\\[8pt]
\twobox{1}{9}      & & & 25     & & &  105 to 145 &\\[8pt]
\twobox{1}{9}      & & & 25     & & &  \Z95 to 155 &\\[8pt]
\twobox{1}{9}      & & & 25     & & &  under 125 &\\[8pt]
\twobox{0}{1}      & & & 144   & & & 66 to 78 &\\[8pt]
\twobox{0}{1}      & & & 144   & & & 60 to 78 &\\[8pt]
\twobox{0}{1}      & & & 10,000 & & & 4950 to 5050 &\\[8pt]
\threebox{0}{1}{2} & & & 100    & & & 92 to 108 & \\
\end{tabular}}
\end{center}

\vfill
$\mbox{EV}_{\mbox{\scriptsize sum}} = n\times\mbox{AV}_{\mbox{\scriptsize box}}$\hspace{20pt}
$\mbox{SE}_{\mbox{\scriptsize sum}} = \sqrt{n}\times\mbox{SD}_{\mbox{\scriptsize box}}$
\eject

%\newcommand{\zobox}[2]{\begin{pspicture}(-0.3,0.2)(2,1.3)
%\psframe(-0.1,0.05)(0.5,0.65)  \rput(0.2,0.35){0}\rput[bl](0.53,0.64){\scriptsize #1}
%\psframe(1,0.05)(1.6,0.65)\rput(1.3,0.35){1}\rput[bl](1.64,0.64){\scriptsize #2}
%\psline(-0.3,0.9)(-0.3,-0.2)(2,-0.2)(2,0.9)(-0.3,0.9)
%\end{pspicture}}

\newcommand{\zobox}[2]{\begin{pspicture}(-0.2,0.2)(2.4,1.3)
\rput[r](-0.08,0.35){\small #1\vphantom{,'}}\psframe(0.0,0.05)(0.5,0.65)\rput(0.25,0.35){0}
   \rput[l](0.58,0.35){s\vphantom{',1}}
\rput[r](1.74,0.35){\small #2\vphantom{,'}}\psframe(1.8,0.05)(2.3,0.65)\rput(2.05,0.35){1}
   \rput[l](2.38,0.35){s\vphantom{',1}}
%\psframe(1,0.05)(1.6,0.65)\rput(1.3,0.35){1}\rput[bl](1.64,0.64){\scriptsize #2}
\psline(-0.9,0.9)(-0.9,-0.2)(2.8,-0.2)(2.8,0.9)(-0.9,0.9)
\end{pspicture}}

\pagestyle{empty}
\lstset{language=R, showspaces=false, showstringspaces=false}

\href{http://www.su.edu}{\includegraphics[height=1.75cm]{sulogo.eps}}
\vspace{-1.79cm}

{{\ }\hfill\small
\begin{tabular}{cl}
& Math 207\\
& Introduction to Statistics\\
\end{tabular}
}
\setlength{\baselineskip}{1.05\baselineskip}
\medskip

\begin{center}
\textbf{\large   Expected Values and Standard Errors for Percentages}
\end{center}

1. Fill in the blanks.  Since the boxes contain only zeros and ones, the
100\% $\times$ $($the average of a sample$)$ is the percentage of ones in the sample.
Compute the $z$-scores for the range then use
the normal curve to compute the probability of getting a sample percentage
in that range.
%Use \texttt{R} or the tables in your book as needed.
\vspace{-8pt}

\begin{center}
{\setlength{\tabcolsep}{8pt}
\begin{tabular}{cccccccc}
\textbf{Box} & $\mbox{mean}_{\mbox{\footnotesize box}}$ & 
   $\mbox{SD}_{\mbox{\footnotesize box}}$ & $n$ & $\mbox{EV}_{\mbox{\scriptsize \%}}$ &
   $\mbox{SE}_{\mbox{\scriptsize \%}}$ & \textbf{range} & \textbf{probability}\\[3pt]\hline
\zobox{5}{5}      & $0.5$ & $0.5$ & 16       & 50\% & 12.5\%& 37.5\% to 62.5\% & 68\%\\[8pt]
\zobox{5}{5}      & & & 16       & & & 25.0\% to   75.0\% &\\[8pt]
\zobox{5}{5}      & & & 16       & & & 45.0\% to   55.0\% &\\[8pt]
\zobox{5}{5}      & & & 100      & & & 37.5\% to 62.5\% &\\[8pt]
\zobox{5}{5}      & & & 100      & & & 25.0\% to   75.0\% &\\[8pt]
\zobox{5}{5}      & & & 100      & & & 45.0\% to   55.0\% &\\[8pt]
\zobox{65}{35}    & & & 100    & & & 30.0\% to   40.0\% &\\[8pt]
\zobox{65}{35}    & & & 1,000  & & & 30.0\% to   40.0\% &\\[8pt]
\zobox{65}{35}    & & & 10,000  & & & 30.0\% to   40.0\% &\\[8pt]
\zobox{3,581}{3,091}    & & & 25  & & & 41.0\% to   51\% &\\[8pt]
\zobox{3,581}{3,091}    & & & 100  & & & 41.0\% to   51\% &\\[8pt]
\zobox{3,581}{3,091}    & & & 1,000  & & & 41.0\% to   51\% &\\[8pt]
\zobox{3,581}{3,091}    & & & 1,000  & & & 45.0\% to   47\% &\\[8pt]
\end{tabular}}
\end{center}


\vfill
$\mbox{EV}_{\mbox{\scriptsize av}} = \mbox{AV}_{\mbox{\scriptsize box}}$\hspace{10pt}
$\mbox{SE}_{\mbox{\scriptsize av}} = \mbox{SD}_{\mbox{\scriptsize box}}\,/\sqrt{n}$
\eject

\end{document}

